\section{Base de Datos: Diseño Conceptual y Estructura Base}
\label{sec:s0_database}

El diseño de la persistencia parte de una premisa heredada del nivel de contexto
(Sección~\ref{ssec:s0_c4_contexto}): la \textbf{autenticación es responsabilidad de Supabase
Auth}, que mantiene su propio schema \comando{auth} con la tabla de usuarios y sus credenciales
cifradas. El dominio de negocio de EthosHub vive en el schema \comando{core}. A diferencia de un
modelo de tabla única, el perfil se \textbf{especializa por rol en dos tablas}:
\comando{core.profiles\_basic} (profesional) y \comando{core.profiles\_company} (reclutador),
ambas vinculadas uno-a-uno con el usuario de Supabase mediante la clave \comando{id\_auth}
(UUID = \comando{auth.users.id}).\par

\subsection{Modelado Entidad-Relación Preliminar}
\label{ssec:s0_er_preliminar}

El modelo conceptual del Sprint 0 fija las entidades fundacionales y deja preparadas las
relaciones que los sprints posteriores poblarán (proyectos, experiencia, formación,
\textit{skills}, conexiones, chat, pipeline de reclutamiento). La Figura~\ref{fig:s0_er} muestra
el núcleo cuenta--perfil y sus extensiones previstas.

\begin{figure}[!ht]
    \centering
    \begin{tikzpicture}[
        ent/.style={draw=colorPrimario, fill=headerTabla, thick, rounded corners=2pt,
            text width=2.7cm, align=center, minimum height=0.9cm, font=\sffamily\scriptsize},
        entx/.style={draw=grisISO, fill=grisTecnico, thick, rounded corners=2pt,
            text width=2.4cm, align=center, minimum height=0.8cm, font=\sffamily\scriptsize},
        rel/.style={-{Latex[length=2mm]}, grisISO, font=\sffamily\tiny}
    ]
        \node[ent] (auth) {auth.users\\\tiny(Supabase)};
        \node[ent, above right=0.4cm and 2.0cm of auth] (basic) {core.profiles\_basic\\\tiny(profesional)};
        \node[ent, below right=0.4cm and 2.0cm of auth] (comp) {core.profiles\_company\\\tiny(reclutador)};
        \node[entx, above right=0.2cm and 1.6cm of basic] (exp) {work\_experiences};
        \node[entx, right=1.6cm of basic] (sk) {hard/soft\_skills};
        \node[entx, below right=0.2cm and 1.6cm of basic] (set) {portfolio\_settings};

        \draw[rel] (auth) -- node[above]{1:1} (basic);
        \draw[rel] (auth) -- node[below]{1:1} (comp);
        \draw[rel] (basic) -- node[above]{1:N} (exp);
        \draw[rel] (basic) -- node[above]{1:N} (sk);
        \draw[rel] (basic) -- node[below]{1:1} (set);
    \end{tikzpicture}
    \caption{Modelo Entidad-Relación preliminar: núcleo cuenta--perfil (dos roles) y extensiones.}
    \label{fig:s0_er}
\end{figure}

\begin{NotaTecnica}
Las entidades en color tenue (\comando{work\_experiences}, \comando{hard\_skills}/\comando{soft\_skills},
\comando{portfolio\_settings}, además de \comando{projects}, \comando{profile\_connections},
\comando{chat\_messages} y \comando{profile\_likes}) se diseñan conceptualmente en el Sprint 0
pero se \textbf{implementan} en los Sprints 1--4 mediante las migraciones \comando{V1}--\comando{V15},
documentadas en el Anexo~\ref{cap:anexo_migraciones}.
\end{NotaTecnica}

\subsection{Entidades Centrales y Esquema Fundacional de la Cuenta del Perfil}
\label{ssec:s0_entidades_centrales}

La tabla \comando{core.profiles\_basic} es el esquema fundacional sobre el que se apoya la
gestión de la cuenta del profesional. Sus atributos base se acordaron en el Sprint 0 y se resumen
en la tabla~\ref{tab:s0_profiles}; el detalle completo (y la tabla análoga
\comando{profiles\_company}) figura en el Anexo~\ref{cap:anexo_datos}.

\begin{table}[!ht]
    \centering
    \renewcommand{\arraystretch}{1.3}
    \fontsize{9}{10.8}\selectfont\sffamily
    \begin{tabularx}{\textwidth}{|l|l|X|}
        \hline
        \rowcolor{colorPrimario}
        \textcolor{white}{\textbf{Campo}} & \textcolor{white}{\textbf{Tipo}} & \textcolor{white}{\textbf{Descripción}} \\ \hline
        \comando{id\_auth} & \comando{uuid} (PK) & Identidad; referencia 1:1 a \comando{auth.users.id}. \\ \hline
        \rowcolor{grisTecnico}
        \comando{first\_name} / \comando{last\_name} & \comando{text} & Nombre y apellidos del profesional. \\ \hline
        \comando{email} & \comando{text} & Correo, sincronizado desde \comando{auth} vía \textit{trigger}. \\ \hline
        \rowcolor{grisTecnico}
        \comando{bio} & \comando{text} & Biografía profesional. \\ \hline
        \comando{avatar\_url} & \comando{text} & URL de la imagen de perfil. \\ \hline
        \rowcolor{grisTecnico}
        \comando{is\_active} & \comando{boolean} & Estado de la cuenta. \\ \hline
        \comando{deleted\_at} & \comando{timestamptz} & Baja lógica (soft-delete). \\ \hline
        \rowcolor{grisTecnico}
        \comando{created\_at} / \comando{updated\_at} & \comando{timestamptz} & Marcas temporales (UTC). \\ \hline
    \end{tabularx}
    \caption{Esquema fundacional de \comando{core.profiles\_basic} (Sprint 0).}
    \label{tab:s0_profiles}
\end{table}

El \textbf{rol} del usuario no es una columna: se deriva de la tabla en la que existe el perfil
(\comando{profiles\_basic} = profesional, \comando{profiles\_company} = reclutador); las cuentas
administrativas se identifican por dominio de correo. Sobre este esquema se fijaron tres
principios de diseño que rigen el resto de la base de datos:

\begin{enumerate}[noitemsep]
    \item \textbf{Aislamiento por RLS:} cada fila de \comando{core} se protege con
    \textit{Row Level Security} de PostgreSQL, de modo que un usuario solo accede a sus
    propios datos.
    \item \textbf{Operaciones sensibles vía RPC:} las escrituras que cruzan la frontera de RLS
    se ejecutan mediante funciones \comando{SECURITY DEFINER} con \comando{search\_path} fijado,
    nunca demoviendo el rol dentro de la función.
    \item \textbf{Marcas temporales en UTC:} todos los \comando{timestamptz} se almacenan en
    UTC para garantizar consistencia entre zonas horarias.
\end{enumerate}

Las convenciones de migración se establecen también aquí: los \textit{scripts} versionados
residen en \comando{src/main/resources/db/migration/} con el formato estilo Flyway
\comando{V<n>__descripcion.sql}. El Sprint 0 entrega el \textit{baseline} del schema
\comando{core}; las migraciones \comando{V1} en adelante se documentan en los capítulos
correspondientes y en el Anexo~\ref{cap:anexo_migraciones}.

\clearpage
